%%% This is a LaTeX template for the small group project self-assessment document.
%%% To compile this document, make sure pdflatex is installed on your workstation and use:
%%% pdflatex self-assessment.tex in a UNIX CLI.
\documentclass[10pt,a4paper]{article}
\usepackage[margin=2.5cm]{geometry}
\usepackage{longtable}
\renewcommand{\familydefault}{\sfdefault}

\begin{document}
\title{Small Group Project Self-Assessment}
\author{}
\date{}
\maketitle

%%% Lines starting with %%% are comment lines.
%%% Fill out the form by removing all comment lines and replacing them by the suggested content.

\noindent\textbf{Team name}: Flax-6

\noindent\textbf{URL of the deployed application}: https://hamed54.pythonanywhere.com/

\noindent\textbf{Team members}:
\begin{itemize}
\setlength\itemsep{0em}
\item Arvin Hadidi-Fard
\item Cate Tay Chew
\item Hamed Latif
\item Katherine Wanting Zhao
\item Kausar Khan
\end{itemize}

\noindent\textbf{If the team agrees on an effective team size, specify it here.  Otherwise, state that the team does not agree.} 

\vspace{2mm}

\hspace{-7mm} The team agrees that the effective team size was five.


%%% Submisison Checklist 
\section*{Submission checklist}
\begin{tabular}{|p{15mm}|p{13cm}|}
\hline
Check & Task \\
\hline
& The work that is submitted has been produced by members of the team only, without input from other students or third parties.  All use of examples from the internet and generative AI are declared in the \texttt{README.md} file.\\
\hline
& The team is familiar with the required deliverables for this project, outlined in the small group project handbook, available on KEATS.\\
\hline
& The team's shared GitHub repository has been registered on Team Feedback.  All commits are attributed to a member of the team (check this on Team Feedback!).  All collaborative coding sessions have been registered by the committer on Team Feedback.  The shared repository has a default branch that corresponds to the accepted version of the software.\\
\hline
& The version of the web application that the team wishes to submit has been deployed.  The URL is recorded in the \texttt{README.md} file.  The deployed system is seeded with data that meets the requirements outlined in the specifications.  The system is accessible using the required access credentials.  The deployed system has been manually tested to ensure all functionality available in development works and is accessible in production.\\
\hline
& All team members listed above agree with the claims made in this checklist.\\
\hline
\end{tabular}


%%% Business Objectives
\section*{Objective}
\textbf{Briefly outline the business objective of the system you have built.  
Aim to identify what the system you developed aims to achieve 
(do not regurgitate the overall assignment).  
This section should consist of a single paragraph.  Aim for less than 100 words.}

\vspace{2mm}

\hspace{-7mm} This project aims to deliver a functional recipe management platform 
for students in shared accommodation who struggle with scattered recipe sources
and meal planning. 
Our system integrates three core workflows: 
users can create and discover recipes with tags, 
plan weekly meals with a reusable calendar interface, 
and generate categorised shopping lists automatically from their meal plans. 
We incorporated AI-powered recipe generation to reduce decision fatigue 
and added social features, allowing users to follow others and 
discover popular recipes. 
The platform combines practical cooking with community features, 
making consistent home cooking more achievable for students 
managing budgets and busy schedules.


\section*{Use cases}
Use cases clarify who uses the system and how they can use it to accomplish particular tasks.  Identify the main use cases of your team's application.  For each use case, provide:
\begin{itemize}
\item Name: a name for the use case.  If the name does not describe the task adequately by itself, add a short task description.
\item Actor \& trigger: Identify who and what starts the interaction with the system.
\item Main success scenario \& outcome: List the main steps that lead to a successful outcome.
\item Alternate flows: Identify possible alternative flow conditions (e.g. what happens under anomalous conditions).
\end{itemize}
An example of a use case is as follows:
\begin{itemize}
\item Name: Place order
\item Actor and trigger: Customer selects products and proceeds to checkout.
\item Main success scenario: (1) Customer logs in, (2) adds products to cart, (3) proceeds to checkout, (4) enters shipping and payment details, (5) confirms order, (6) system processes payment and displays confirmation. Outcome: new order is placed.
\item Alternate flows: Payment fails $\rightarrow$ system redirects to step (4) with error message.
\end{itemize}

\begin{longtable}{|p{30mm}|p{30mm}|p{40mm}|p{40mm}|}
\hline
\textbf{Name (\& short task description)}
& 
\textbf{Actor \& trigger}
& 
\textbf{Main success scenario \& outcome}
& 
\textbf{Alternative flows}\\


%%% Above is the structure to follow %%%

%%% Use Case 1 

\hline
Create Recipe 
& 
Logged-in user selects ``Create Recipe'' from the navigation bar 
&
(1) System displays Step 1 form, 

\vspace{1 mm}

(2) User enters title, publication date, servings, ingredients, steps, optional image, visibility and times.

\vspace{1 mm}

(3) User clicks ``Next: Add Tags'',

\vspace{1 mm}

(4) System displays Step 2 tag page, 

\vspace{1 mm}

(5) User selects tags or adds custom tags,

\vspace{1 mm}

(6) User clicks ``Create Recipe''.  
Outcome: recipe is saved and shown on the detail page with a success message. 
&
Missing fields $\rightarrow$ system shows validation errors. 

\vspace{1 mm}

Invalid date $\rightarrow$  user must enter today or earlier.

\vspace{1 mm}

Image upload fails user may retry or continue without image.  

\vspace{1 mm}

User cancels $\rightarrow$  no recipe is saved and user is taken back to main page. \\
\hline


%%% Use Case 2
\hline
Filter Recipe
&
Logged-in user opens the all recipes page by clicking "View All" on the dashboard page  
or by clicking "Discover $\rightarrow$ All Recipes" in the navigation bar.
&
(1) System displays all public recipes (default: newest first).

\vspace{1 mm}

(2) User selects one or more tag filters (Course, Dietary, Cuisine, Others).

\vspace{1 mm}

(3) User applies filters by pressing Return or clicking "Apply".  

Outcome: Active filters will be shown above all the recipes and below the filter buttons.

\vspace{1 mm}

(4) User clears filters to return to the full recipe list.

Outcome: All recipes will be shown again

\vspace{1 mm}

(5) User can add to the filters by clicking (checking) more filters

Outcome: Recipes shown gets further filtered to show recipes that contain all tag(s) selected
&
No alternate flows. \\


%%% Use Case 3
\hline
Sort Recipes
&
Logged-in user opens the all recipes page by clicking view all on the dashboard page or clicking Discover/All Recipes in the navigation bar.
&
(1) Users click the the drop down that is default "Newest First".

\vspace{1 mm}

(2) From the dropdown options (publication date, title, time, most saved, highest rated) the user will click what they want to sort the recipes by.

Outcome: System displays recipes sorted based on the sort by option selected.
&
e.g. if "Prep Time (Least)" was selected, the recipes shown will start from the recipe with the shortest prep time.

\vspace{1 mm}

(3) If a filter selected before or after sorting, the filtered recipes will be sorted. \\


%%% Use Case 4
\hline
Update my own recipe
&
Logged-in User opens My Kitchen/My Recipes in the navigation bar to view all their own recipes. 

User clicks edit for the recipe they want to Update
&
(1) System displays the first page (Step 1) of the same form as the create recipe form, 
but with the fields already filled in with the existing recipe's details.

\vspace{1 mm}

(2) User can update any field they want by simply deleting and/or adding words/numbers.

If there's an image present, it will be shown at the bottom of the page with a checkbox "Delete image".

\vspace{1 mm}

(3) Click the delete image checkbox to delete the image.

\vspace{1 mm}

(4) To replace the existing image, add a new image under the image field.

\vspace{1 mm}

(5) Click "Next: Add Tags" 

Outcome: System displays "Step 2" tags page. 

\vspace{1 mm}

(6) In Step 2, the user can delete/add tags.

\vspace{1 mm}

(7) User clicks "Update Recipe"

Outcome: System displays the recipe detail page for the updated Recipe.

At the top of the page it will show that the recipe was successfully updated.
&
Missing fields $\rightarrow$ system shows validation errors. 

\vspace{1 mm}

Invalid date $\rightarrow$  user must enter today or earlier.

\vspace{1 mm}

Image upload fails user may retry or continue without image.  

\vspace{1 mm}

User cancels $\rightarrow$  recipe remains the same with no updates. \\


%%% Use Case 5
\hline
Create a modification of another user's recipe
&
Logged-in user views a recipe by clicking on it. 

If the recipe doesn't belong to them, they can create a modication by clicking "Create My Version" at the bottom of the page.
&
(1) System displays the "Step 1" of the create recipe form, but with the fields filled in with the existing recipe's details.

At the top of the form, it will tell the user that it is creating a modified version 
of the original author's (showing the actual username) recipe.

\vspace{1 mm}

(2) User can add/delete to the fields. 

If there's an image present, it will be shown at the bottom of the page with a checkbox "Delete image".

\vspace{1 mm}

(3) Click the delete image checkbox to delete the image.

\vspace{1 mm}

(4) To replace the existing image, add a new image under the image field.

\vspace{1 mm}

(5) Click "Next: Add Tags" 

Outcome: System displays "Step 2" tags page. 

\vspace{1 mm}

(6) In Step 2, the user can delete/add tags.

\vspace{1 mm}

(7) User clicks "Create My Version"

Outcome: System displays the recipe detail page for the new Recipe created.

At the top of the page it will show that the recipe was successfully created (authored by the user).

The original recipe will remain untouched.
&
Missing fields $\rightarrow$ system shows validation errors. 

\vspace{1 mm}

Invalid date $\rightarrow$  user must enter today or earlier.

\vspace{1 mm}

Image upload fails user may retry or continue without image.  

\vspace{1 mm}

User cancels $\rightarrow$  original recipe remains unchanged. \\


%%% Use case 6
\hline
Browse and Search Users
&
Logged-in user clicks "View All" in dashboard page 
or Discover/All Recipes in navigation bar.

System displays the all recipes page.

User clicks the user tab 
&
(1) System displays all users

\vspace{1 mm}

(2) User clicks follow/unfollow.

Outcome: System displays newly followed/unfollowed user.

System tells the user at the top of the page that they followed/unfollowed the user shown.

\vspace{1 mm}

(3) In the view of all users, the user can view the profile by clicking the user card

Outcome: System displays the profile of the user clicked

\vspace{1 mm}

(4) User can search by username (@johndoe), first name (john), last name (doe), or both (john doe)

Outcome: System displays recipes made by users matching the search query.
&
User can click user tab to view users that match the search query.

\vspace{1 mm}

If no user/recipe matches the search query, no search results will be shown. \\
\hline
\end{longtable}

\end{document}